\documentclass[11pt,a4paper]{article}

\usepackage[spanish,es-noshorthands]{babel}
\usepackage[margin=2.5cm]{geometry}
\usepackage{xcolor}
\usepackage{listings}
\usepackage{booktabs}
\usepackage{array}
\usepackage{enumitem}
\usepackage{forest}
\usepackage{tikz}
\usetikzlibrary{positioning,arrows.meta}
\usepackage[hidelinks]{hyperref}

\definecolor{gris}{HTML}{6B7280}
\definecolor{fondo}{HTML}{F6F7F9}
\definecolor{borde}{HTML}{E3E6EA}
\definecolor{acento}{HTML}{2F6B4F}

\lstdefinestyle{codigo}{
  basicstyle=\ttfamily\footnotesize,
  backgroundcolor=\color{fondo},
  frame=single,
  rulecolor=\color{borde},
  framesep=6pt,
  breaklines=true,
  showstringspaces=false,
  keywordstyle=\color{acento}\bfseries,
  commentstyle=\color{gris}\itshape,
  columns=fixed,
  keepspaces=true,
}
\lstdefinelanguage{Elixir}{
  morekeywords={def,defp,defmodule,defstruct,do,end,fn,if,else,case,cond,true,false,nil,for,reduce,into,when},
  morecomment=[l]{\#},
  morestring=[b]",
}
\lstset{style=codigo}

\newcommand{\gram}[1]{\texttt{#1}}

\title{\vspace{-1.5cm}\textbf{Proyecto 2 — Algoritmo CYK}\\[0.3cm]
\large Simplificación de gramáticas libres de contexto y parsing}

\author{%
  Harry Mendez (24089) \\
  Juan Gualim (24852) \\
  Diego Parada (2020225)
}

\date{Teoría de la Computación \\ Universidad del Valle de Guatemala \\[0.2cm] Octubre 2026}

\begin{document}
\maketitle

\begin{center}
\small Repositorio: \url{https://github.com/hmndz3/Proyecto2TC}
\end{center}

\vspace{0.5cm}
\tableofcontents
\newpage

% ====================================================================
\section{Introducción}

Este proyecto responde a una pregunta sencilla de enunciar y difícil de
resolver: \textbf{dada una gramática y una frase, ¿esa frase pertenece al
lenguaje que la gramática genera?}

Para una gramática pequeña uno puede intentar derivar la frase a mano, como
se hizo en los laboratorios del curso. Pero ese método tiene un problema: si
la gramática es ambigua o tiene recursión, uno puede quedarse probando
derivaciones para siempre sin llegar a una respuesta. El algoritmo
\textbf{CYK} (Cocke--Younger--Kasami) resuelve esto: siempre termina, siempre
da una respuesta correcta, y lo hace en tiempo polinomial.

El precio que cobra CYK es que la gramática tiene que estar en
\textbf{Forma Normal de Chomsky} (CNF): toda producción debe ser de la forma
$A \to BC$ o $A \to a$. Casi ninguna gramática viene así escrita, por eso la
primera mitad del proyecto es convertirla, y esa conversión tiene sus propias
reglas que hay que aplicar en orden.

Entonces el programa hace dos cosas:

\begin{enumerate}[leftmargin=*]
  \item \textbf{Simplifica} la gramática y la lleva a CNF, en cuatro pasos.
  \item \textbf{Valida} una frase con CYK, reportando sí o no, cuánto tardó y
  el árbol de derivación.
\end{enumerate}

Está escrito en \textbf{Elixir}, sin librerías externas, y se usa desde la
terminal.

% ====================================================================
\section{Marco: gramáticas y lenguajes}

\subsection{Qué es una gramática}

Una gramática libre de contexto se describe con cuatro cosas,
$G = (V, T, P, S)$:

\begin{itemize}[leftmargin=*]
  \item $V$: las \textbf{variables} o símbolos no terminales. Son los que se
  pueden seguir reemplazando. Por convención van en mayúscula: \gram{S},
  \gram{NP}, \gram{VP}.
  \item $T$: los \textbf{terminales}. Son las palabras reales del lenguaje,
  las que ya no se reemplazan por nada: \gram{she}, \gram{eats}, \gram{id},
  \gram{+}.
  \item $P$: las \textbf{producciones} o reglas, del tipo \gram{A -> B C}.
  \item $S$: el \textbf{símbolo inicial}, por donde empieza toda derivación.
\end{itemize}

«Libre de contexto» significa que el lado izquierdo de cada regla es siempre
una sola variable. No existe algo como \gram{a B -> c}: la variable \gram{B}
se reemplaza igual sin importar qué tenga a los lados. Eso es lo que la
distingue de las gramáticas sensibles al contexto en la jerarquía de Chomsky,
y lo que hace que exista un algoritmo eficiente para reconocerlas.

\subsection{La gramática del enunciado}

\begin{lstlisting}
S   -> NP VP
VP  -> VP PP | V NP | cooks | drinks | eats | cuts
PP  -> P NP
NP  -> Det N | he | she
V   -> cooks | drinks | eats | cuts
P   -> in | with
N   -> cat | dog | beer | cake | juice | meat | soup | fork | knife | oven | spoon
Det -> a | the
\end{lstlisting}

Leída en voz alta: una oración (\gram{S}) es un sujeto (\gram{NP}) seguido de
un predicado (\gram{VP}). Un \gram{NP} es un determinante más un sustantivo
(\gram{the cat}) o un pronombre (\gram{he}, \gram{she}). Un \gram{VP} es un
verbo con su objeto (\gram{eats a cake}), o un \gram{VP} con una frase
preposicional pegada (\gram{eats a cake with a fork}), o un verbo solo
(\gram{cooks}).

\subsection{Por qué hace falta la Forma Normal de Chomsky}

CYK funciona partiendo la frase en dos pedazos y preguntando, para cada
partición, qué variable puede generar el pedazo izquierdo y cuál el derecho.
Eso solo tiene sentido si las reglas tienen exactamente dos símbolos del lado
derecho. Una regla como \gram{S -> A B C d} no dice dónde partir.

Por eso toda producción debe quedar de una de estas dos formas:

\begin{center}
\begin{tabular}{ll}
\toprule
$A \to B\,C$ & dos variables, ninguna terminal \\
$A \to a$    & una sola terminal \\
\bottomrule
\end{tabular}
\end{center}

Y además no puede haber producciones vacías ($A \to \varepsilon$) ni unitarias
($A \to B$), porque romperían el conteo de longitudes en que se basa el
algoritmo.

% ====================================================================
\section{Diseño de la aplicación}

\subsection{Por qué Elixir}

El enunciado no pide ningún lenguaje en particular. Elegimos Elixir por tres
razones concretas:

\begin{enumerate}[leftmargin=*]
  \item \textbf{Los tres pasos de simplificación son cálculos de punto fijo.}
  Todos tienen la misma forma: se aplica una regla sobre un conjunto una y otra
  vez, hasta que el conjunto deja de crecer. En un lenguaje funcional eso se
  escribe como una función recursiva corta que compara el conjunto anterior con
  el nuevo y se detiene cuando son iguales. No hacen falta variables mutables
  ni banderas de control.
  \item \textbf{Los conjuntos y los mapas vienen con el lenguaje.} Elixir trae
  \texttt{MapSet} y \texttt{Map} listos para usar, y resultan ser justo las dos
  estructuras que este problema necesita: conjuntos para los anulables, los
  generadores y los alcanzables; mapas para las reglas de la gramática y para
  la tabla de CYK. Operaciones como «agregar al conjunto si todavía no está» o
  «buscar todas las reglas que tengan este par» quedan en una línea, sin tener
  que programar la estructura de datos.
  \item \textbf{No necesita instalación de dependencias.} Todo el proyecto usa
  solo la librería estándar.
\end{enumerate}

\subsection{Organización en módulos}

La aplicación está separada en dos capas. Los módulos de abajo son funciones
puras: reciben una gramática y devuelven otra, sin tocar la pantalla ni
archivos. Los de arriba solo se encargan de mostrar.

\begin{center}
\begin{tabular}{@{}>{\ttfamily}l p{9.2cm}@{}}
\toprule
\textrm{\textbf{Archivo}} & \textbf{Responsabilidad} \\
\midrule
\multicolumn{2}{@{}l}{\textit{Lógica (funciones puras)}} \\
lib/cyk/grammar.ex  & Lee el formato \texttt{.txt}, guarda la gramática, la imprime de vuelta. \\
lib/cyk/simplify.ex & Pasos 1 a 3: épsilon, unitarias, inútiles. \\
lib/cyk/cnf.ex      & Paso 4: Forma Normal de Chomsky. \\
lib/cyk/parser.ex   & Algoritmo CYK, tabla dinámica, medición de tiempo. \\
lib/cyk/tree.ex     & Reconstruye el árbol desde los backpointers. \\
lib/cyk.ex          & Encadena los cuatro pasos y parte la frase en tokens. \\
\midrule
\multicolumn{2}{@{}l}{\textit{Presentación}} \\
lib/cyk/ui.ex       & Formato de la salida en terminal: cajas, colores, tablas. \\
lib/cyk/svg.ex      & Dibuja el árbol y lo abre en una ventana aparte. \\
lib/cyk/cli.ex      & Interfaz de línea de comandos y modo interactivo. \\
lib/mix/tasks/cyk.ex & Registra el comando \texttt{mix cyk}. \\
\bottomrule
\end{tabular}
\end{center}

\subsection{Flujo del programa}

\begin{center}
\begin{tikzpicture}[
  caja/.style={draw=borde, fill=fondo, rounded corners=3pt, minimum height=9mm,
               inner xsep=5pt, font=\small},
  flecha/.style={-latex, draw=gris, thick}
]
\node[caja] (a) {\texttt{.txt}};
\node[caja, right=5mm of a] (b) {Gramática};
\node[caja, right=5mm of b] (c) {Paso 1--3};
\node[caja, right=5mm of c] (d) {Paso 4: CNF};
\node[caja, below=7mm of d] (e) {Tabla CYK};
\node[caja, left=5mm of e] (f) {Backpointers};
\node[caja, left=5mm of f] (g) {Árbol};
\node[caja, left=5mm of g] (h) {Sí / No};

\draw[flecha] (a) -- (b);
\draw[flecha] (b) -- (c);
\draw[flecha] (c) -- (d);
\draw[flecha] (d) -- (e);
\draw[flecha] (e) -- (f);
\draw[flecha] (f) -- (g);
\draw[flecha] (e.south) -- ++(0,-6mm) -| (h.south);
\end{tikzpicture}
\end{center}

\subsection{Formato de entrada}

Mantuvimos el mismo formato de los archivos que entregó el curso
(\texttt{1.txt} y \texttt{1-cnf.txt}), para que el programa pueda leerlos tal
cual:

\begin{lstlisting}
# los comentarios empiezan con numeral
S   -> NP VP
VP  -> VP PP | V NP | eats
X   -> + T X | e
\end{lstlisting}

Las reglas del formato son:

\begin{itemize}[leftmargin=*]
  \item \gram{->} separa la cabeza del cuerpo; \gram{|} separa alternativas.
  \item \textbf{Los símbolos van separados por espacios.} Por eso \gram{id} es
  un solo token y \gram{( E )} son tres. Lo mismo al escribir la frase:
  \gram{( id + id ) * id}.
  \item \gram{e} representa épsilon.
  \item Un símbolo es \textbf{no terminal} si aparece como cabeza de alguna
  producción. Todo lo demás es terminal. Esto es más robusto que fijarse en
  mayúsculas y permite nombres como \gram{A'}.
  \item El \textbf{símbolo inicial} es la cabeza de la primera producción, tal
  como lo define la presentación del curso.
\end{itemize}

Las gramáticas incluidas están en la carpeta \texttt{gramaticas/}:

\begin{center}
\begin{tabular}{@{}>{\ttfamily}l p{8.5cm}@{}}
\toprule
\textrm{\textbf{Archivo}} & \textbf{Contenido} \\
\midrule
ingles.txt          & La del enunciado del proyecto. \\
aritmetica.txt      & Expresiones aritméticas: el archivo \texttt{1.txt} del curso. \\
aritmetica-cnf.txt  & El CNF de referencia (\texttt{1-cnf.txt}), usado en las pruebas. \\
ambigua.txt         & \texttt{E -> E + E | E * E | ( E ) | id}, para ver el conteo de árboles. \\
\bottomrule
\end{tabular}
\end{center}

% ====================================================================
\section{El algoritmo de simplificación}

Los cuatro pasos se aplican en este orden, que es el que se vio en clase.
El orden importa: por ejemplo, eliminar épsilon \emph{crea} producciones
unitarias nuevas, así que el paso 2 tiene que ir después del paso 1.

Para ilustrarlos usamos la gramática aritmética, porque la del enunciado
—como se explica más adelante— ya viene en CNF y no ejercita ningún paso.

\begin{lstlisting}
E -> T X
X -> + T X | e
T -> F Y
Y -> * F Y | e
F -> ( E ) | id
\end{lstlisting}

\subsection{Paso 1: eliminación de producciones épsilon}

\textbf{Qué busca.} Las producciones que generan la cadena vacía. Estorban
porque CYK razona sobre longitudes de subcadena, y una variable que puede
desaparecer rompe esa cuenta.

\textbf{Cómo funciona.} Primero se calcula el conjunto de \textbf{anulables}:
toda variable que tenga una producción vacía, más toda variable cuyo cuerpo
esté formado únicamente por anulables. Se repite hasta que el conjunto no
crezca. Después, por cada cuerpo, se generan todas las versiones posibles
quitando o dejando cada símbolo anulable, y se descartan los cuerpos que
quedaron vacíos.

\textbf{En este ejemplo.} Anulables: \gram{X}, \gram{Y}.

\begin{lstlisting}
E -> T X | T
X -> + T X | + T
T -> F Y | F
Y -> * F Y | * F
F -> ( E ) | id          8 -> 10 producciones
\end{lstlisting}

Un detalle: si el símbolo inicial resultaba anulable, significa que la cadena
vacía sí pertenece al lenguaje. Como esa información se perdería al borrar las
épsilon, el programa la guarda aparte en una bandera de la gramática.

\textbf{Dónde está.} \texttt{lib/cyk/simplify.ex}, funciones
\texttt{epsilon/1} y \texttt{nullable/1}.

\subsection{Paso 2: eliminación de producciones unitarias}

\textbf{Qué busca.} Las reglas del tipo \gram{A -> B}, donde una variable
simplemente se renombra como otra. No aportan estructura y CNF no las permite.

\textbf{Cómo funciona.} Se arma el conjunto de \textbf{pares unitarios}
$(A, B)$ tales que $A \Rightarrow^* B$ usando solo reglas unitarias. Se
arranca con los reflexivos $(A, A)$ y se cierra transitivamente. Luego, a cada
$A$ se le copian las producciones \textbf{no} unitarias de cada $B$ de su par.

\textbf{En este ejemplo.} Pares: \gram{(E,F)}, \gram{(E,T)}, \gram{(T,F)}
además de los reflexivos.

\begin{lstlisting}
E -> T X | ( E ) | id | F Y
X -> + T X | + T
T -> ( E ) | id | F Y
Y -> * F Y | * F
F -> ( E ) | id          10 -> 13 producciones
\end{lstlisting}

\textbf{Dónde está.} \texttt{lib/cyk/simplify.ex}, funciones
\texttt{unit/1} y \texttt{unit\_pairs/1}.

\subsection{Paso 3: eliminación de símbolos inútiles}

\textbf{Qué busca.} Variables que no sirven para nada. Son de dos tipos y se
eliminan en ese orden:

\begin{enumerate}[leftmargin=*]
  \item \textbf{No generadores}: variables que nunca llegan a producir una
  cadena de puros terminales. Por ejemplo \gram{C -> a C b} sin otra
  alternativa: siempre queda una \gram{C} pendiente.
  \item \textbf{No alcanzables}: variables a las que nunca se llega desde el
  símbolo inicial.
\end{enumerate}

El orden importa. Si se quitan primero los no alcanzables y después los no
generadores, puede quedar basura: eliminar una variable no generadora puede
dejar a otras sin forma de alcanzarse.

\textbf{En este ejemplo.} No hay ninguno de los dos tipos, la gramática queda
igual (13 producciones). Donde sí se nota es en un caso como el siguiente:

\begin{lstlisting}
Original:                    Despues de los pasos 1 y 2:
  S -> a S | A | C             S -> a S | a | a C b
  A -> a                       A -> a
  B -> a a                     B -> a a
  C -> a C b                   C -> a C b
\end{lstlisting}

\begin{lstlisting}
C no genera terminales      -> se elimina C (y la regla S -> a C b)
A y B ya no son alcanzables -> se eliminan

resultado:  S -> a S | a
\end{lstlisting}

\textbf{Dónde está.} \texttt{lib/cyk/simplify.ex}, funciones
\texttt{useless/1}, \texttt{generating/1} y \texttt{reachable/1}.

\subsection{Paso 4: Forma Normal de Chomsky}

\textbf{Qué busca.} Dejar toda producción como $A \to BC$ o $A \to a$. Son dos
sub-pasos:

\begin{enumerate}[leftmargin=*]
  \item \textbf{Sacar los terminales de los cuerpos largos.} Un cuerpo como
  \gram{( E )} mezcla terminales con variables. Cada terminal se reemplaza por
  una variable nueva que solo produce ese terminal: \gram{T1 -> (}.
  \item \textbf{Binarizar.} Un cuerpo de tres o más símbolos se parte en
  cadena: \gram{A -> B C D} se vuelve \gram{A -> B X1} junto con
  \gram{X1 -> C D}.
\end{enumerate}

\textbf{En este ejemplo.} Variables nuevas: \gram{T1}$\to$\gram{(},
\gram{T2}$\to$\gram{)}, \gram{T3}$\to$\gram{+}, \gram{T4}$\to$\gram{*},
\gram{X1}$\to$\gram{E T2}, \gram{X2}$\to$\gram{T X}, \gram{X3}$\to$\gram{F Y}.

\begin{lstlisting}
E  -> T X | T1 X1 | id | F Y
X  -> T3 X2 | T3 T
T  -> T1 X1 | id | F Y
Y  -> T4 X3 | T4 F
F  -> T1 X1 | id
T1 -> (        T2 -> )        T3 -> +        T4 -> *
X1 -> E T2     X2 -> T X      X3 -> F Y

                                              13 -> 20 producciones
\end{lstlisting}

\textbf{Comparación con el CNF de referencia.} El curso entregó el archivo
\texttt{1-cnf.txt} con la conversión esperada. Nuestro resultado tiene las
mismas 20 producciones y la misma estructura; solo cambian los nombres de las
variables nuevas, que son arbitrarios:

\begin{center}
\begin{tabular}{@{}llllllll@{}}
\toprule
Referencia & \gram{L} & \gram{R} & \gram{P} & \gram{M} & \gram{Z} & \gram{W} & \gram{V} \\
Nuestro    & \gram{T1} & \gram{T2} & \gram{T3} & \gram{T4} & \gram{X1} & \gram{X2} & \gram{X3} \\
\bottomrule
\end{tabular}
\end{center}

\textbf{Dónde está.} \texttt{lib/cyk/cnf.ex}.

% ====================================================================
\section{El algoritmo CYK}

\subsection{La idea}

CYK es programación dinámica. En vez de intentar derivar la frase de arriba
hacia abajo, construye la respuesta de abajo hacia arriba: primero resuelve
las subcadenas de una palabra, después las de dos, y así hasta cubrir la
frase completa.

La pregunta que responde cada celda de la tabla es: \textit{¿qué variables
pueden generar esta subcadena?}

\begin{itemize}[leftmargin=*]
  \item \textbf{Subcadenas de longitud 1.} Para la palabra \gram{cake}, se
  buscan todas las variables con una regla \gram{A -> cake}. Aquí: \gram{N}.
  \item \textbf{Subcadenas más largas.} Se prueban todas las formas de partir
  la subcadena en dos. Para cada partición, si la mitad izquierda la puede
  generar \gram{B}, la derecha la puede generar \gram{C}, y existe una regla
  \gram{A -> B C}, entonces \gram{A} puede generar la subcadena entera.
\end{itemize}

Al final, la frase pertenece al lenguaje si y solo si el \textbf{símbolo
inicial} aparece en la celda que cubre la frase completa.

\subsection{La tabla, con un ejemplo}

Para \textit{she eats a cake with a fork}, con las filas numeradas por
longitud de subcadena:

\begin{center}
\ttfamily\footnotesize
\begin{tabular}{@{}r|ccccccc@{}}
   & she & eats & a   & cake & with & a   & fork \\
\midrule
 1 & NP  & V,VP & Det & N    & P    & Det & N \\
 2 & S   & ·    & NP  & ·    & ·    & NP  &   \\
 3 & ·   & VP   & ·   & ·    & PP   &     &   \\
 4 & S   & ·    & ·   & ·    &      &     &   \\
 5 & ·   & ·    & ·   &      &      &     &   \\
 6 & ·   & VP   &     &      &      &     &   \\
 7 & S   &      &     &      &      &     &   \\
\end{tabular}
\end{center}

Se lee así: la celda de la fila 2, columna \gram{a}, contiene \gram{NP}, lo
que dice que \textit{a cake} puede generarse desde \gram{NP}. La celda de
abajo a la izquierda (fila 7, columna \gram{she}) cubre la frase entera y
contiene \gram{S}, el símbolo inicial. Por lo tanto la frase \textbf{sí}
pertenece al lenguaje.

\subsection{Backpointers y el árbol}

Si cada celda guardara solo «sí, esta variable puede», la respuesta sería un
booleano y nada más. Para poder construir el árbol, cada entrada guarda además
\textbf{de dónde salió}: con qué regla y en qué punto se partió la subcadena.
Eso es un \textit{backpointer}.

Reconstruir el árbol es entonces seguir esos punteros hacia atrás desde la
celda final. El resultado para la frase del ejemplo:

\begin{center}
\begin{forest}
for tree={font=\small\ttfamily, s sep=4mm, l sep=6mm, edge={draw=gris}}
[S
  [NP [she, font=\small\ttfamily\color{acento}]]
  [VP
    [VP
      [V [eats, font=\small\ttfamily\color{acento}]]
      [NP [Det [a, font=\small\ttfamily\color{acento}]] [N [cake, font=\small\ttfamily\color{acento}]]]
    ]
    [PP
      [P [with, font=\small\ttfamily\color{acento}]]
      [NP [Det [a, font=\small\ttfamily\color{acento}]] [N [fork, font=\small\ttfamily\color{acento}]]]
    ]
  ]
]
\end{forest}
\end{center}

El programa muestra este árbol de dos maneras: dibujado en texto dentro de la
terminal, y dibujado como gráfico en una ventana aparte que se abre
automáticamente.

Los backpointers sirven para una cosa más: \textbf{contar cuántos árboles
distintos existen}, que es exactamente la definición de ambigüedad. Si una
celda tiene dos backpointers para la misma variable, hay dos derivaciones.

\subsection{Complejidad}

Son tres bucles anidados —longitud, posición inicial, punto de partición— y
dentro de cada uno se recorren las reglas. Eso da
$\mathcal{O}(n^3 \cdot |G|)$, donde $n$ es la cantidad de palabras y $|G|$ el
tamaño de la gramática.

Para las frases de este proyecto eso es nada: una frase de siete palabras se
resuelve en unos 14 microsegundos.

\textbf{Dónde está.} \texttt{lib/cyk/parser.ex} (tabla y conteo de árboles),
\texttt{lib/cyk/tree.ex} (reconstrucción), \texttt{lib/cyk/svg.ex} (dibujo).

% ====================================================================
\section{Cómo ejecutar el programa}

Lo único que hay que tener instalado es Elixir 1.14 o superior. El proyecto no
usa librerías externas, así que no hay dependencias que descargar.

\begin{lstlisting}
mix cyk                              # modo interactivo
mix cyk gramaticas/aritmetica.txt    # otra gramatica
mix cyk --frase "she eats a cake"    # valida una frase y sale
mix cyk --pasos                      # muestra cada gramatica intermedia
mix cyk --sin-ventana                # no abre el arbol en el navegador
mix cyk --repetir 1000               # fija las corridas de la medicion
mix test                             # las pruebas automaticas
\end{lstlisting}

Dentro del modo interactivo se pueden usar los comandos \gram{:gramatica},
\gram{:cnf}, \gram{:pasos} y \gram{:salir}.

La salida muestra, en orden: el resumen de la gramática, los cuatro pasos de
simplificación con cuántas producciones entran y salen de cada uno, la
gramática en CNF, el veredicto, el tiempo, la cantidad de árboles, la tabla
CYK y el árbol de derivación.

% ====================================================================
\section{Ejemplos y pruebas}

\subsection{Frases que sí están en el lenguaje}

\begin{center}
\begin{tabular}{@{}>{\ttfamily}l p{7.5cm}@{}}
\toprule
\textrm{\textbf{Frase}} & \textbf{Por qué} \\
\midrule
she eats a cake with a fork & \gram{NP VP}, con \gram{VP -> VP PP}. \\
the cat drinks the beer     & Los dos \gram{NP} son \gram{Det N}. \\
he drinks the beer          & \gram{NP} pronombre. \\
she cooks                   & \gram{VP} de un solo verbo. \\
a dog cuts the meat with a knife & El caso más largo, siete palabras. \\
he eats the soup in the oven & Usa la otra preposición. \\
the cake eats               & Correcta sintácticamente aunque absurda. \\
\bottomrule
\end{tabular}
\end{center}

El último caso merece un comentario. \textit{The cake eats} se acepta porque
\gram{the cake} es un \gram{NP} válido y \gram{eats} es un \gram{VP} válido.
La gramática no sabe nada de qué cosas pueden comer. El enunciado lo advierte
explícitamente: las frases aceptadas son sintácticamente correctas, nadie
garantiza que lo sean semánticamente.

\subsection{Frases que no están en el lenguaje}

\begin{center}
\begin{tabular}{@{}>{\ttfamily}l p{7.5cm}@{}}
\toprule
\textrm{\textbf{Frase}} & \textbf{Por qué falla} \\
\midrule
eats she a cake     & Empieza con verbo; \gram{S} exige \gram{NP} primero. \\
a cat with the dog  & Es \gram{NP PP}, le falta el \gram{VP}. \\
she eats a          & \gram{a} es \gram{Det} y queda sin sustantivo. \\
cat the drinks beer & \gram{cat} suelto es \gram{N}, no \gram{NP}. \\
she eats a cake with & La preposición queda colgando sin \gram{NP}. \\
banana              & No es terminal de la gramática. \\
\bottomrule
\end{tabular}
\end{center}

En el último caso el programa además avisa \texttt{Desconocidos: banana}, para
distinguir «la frase está mal armada» de «esa palabra no existe en la
gramática».

\subsection{Expresiones aritméticas}

Con \texttt{gramaticas/aritmetica.txt}. Recordar que los símbolos van
separados por espacios.

\begin{center}
\begin{tabular}{@{}>{\ttfamily}l c >{\ttfamily}l c@{}}
\toprule
\textrm{\textbf{Acepta}} & & \textrm{\textbf{Rechaza}} & \\
\midrule
id                  & sí & + id     & no \\
id + id             & sí & id +     & no \\
( id )              & sí & ( id     & no \\
( id + id ) * id    & sí & id id    & no \\
id * ( id + id )    & sí & * id *   & no \\
\bottomrule
\end{tabular}
\end{center}

\subsection{Ambigüedad}

Con \texttt{gramaticas/ambigua.txt}, que contiene
\gram{E -> E + E | E * E | ( E ) | id}. Esta gramática no define precedencia,
así que una misma expresión admite varias lecturas:

\begin{center}
\begin{tabular}{@{}>{\ttfamily}l c@{}}
\toprule
\textrm{\textbf{Expresión}} & \textbf{Árboles} \\
\midrule
id + id             & 1 \\
id + id * id        & 2 \\
id + id * id + id   & 5 \\
\bottomrule
\end{tabular}
\end{center}

Los números 1, 2, 5 son los números de Catalan, que es justo la cantidad de
formas de parentizar una expresión. El programa reporta esta cuenta junto al
veredicto.

\subsection{Tiempos de ejecución}

Medidos en la máquina de desarrollo, con la gramática del enunciado:

\begin{center}
\begin{tabular}{@{}lcc@{}}
\toprule
\textbf{Frase} & \textbf{Palabras} & \textbf{Tiempo} \\
\midrule
\texttt{she cooks}                   & 2 & 7.9 µs \\
\texttt{cat the drinks beer}         & 4 & 13.4 µs \\
\texttt{she eats a cake with a fork} & 7 & 13.8 µs \\
\bottomrule
\end{tabular}
\end{center}

\subsection{Pruebas automáticas}

El proyecto incluye \textbf{41 pruebas} que corren con \texttt{mix test}.
Están en la carpeta \texttt{test/}:

\begin{center}
\begin{tabular}{@{}>{\ttfamily}l p{8cm}@{}}
\toprule
\textrm{\textbf{Archivo}} & \textbf{Qué verifica} \\
\midrule
grammar\_test.exs  & Lectura del formato, comentarios, épsilon, símbolo inicial. \\
simplify\_test.exs & Los tres pasos, contra gramáticas pequeñas resueltas a mano paso por paso. \\
cnf\_test.exs      & Que toda producción quede $A \to BC$ o $A \to a$, y la equivalencia con el CNF de referencia. \\
cyk\_test.exs      & Frases aceptadas y rechazadas, conteo de árboles, forma del árbol. \\
\bottomrule
\end{tabular}
\end{center}

La prueba de equivalencia merece explicación. No se puede comparar nuestro CNF
contra el \texttt{1-cnf.txt} del curso línea por línea, porque los nombres de
las variables nuevas son arbitrarios (nuestro \gram{T1} es su \gram{L}). Lo
que sí tiene sentido es comparar \textbf{comportamiento}: la prueba corre CYK
con las dos gramáticas sobre una lista de cadenas y verifica que acepten y
rechacen exactamente las mismas.

% ====================================================================
\section{Discusión}

\subsection{Obstáculos encontrados}

\paragraph{Programación dinámica sin variables mutables.}
Este fue el obstáculo conceptual principal. CYK se enseña con una matriz que
se va llenando, y en Elixir no existen las variables mutables. La solución
salió de observar una propiedad del algoritmo: CYK llena la tabla
\textbf{por longitud creciente}, así que cuando se calcula una celda, todas
las que necesita ya están escritas y solamente se leen. Nunca hay que
modificar algo ya calculado. Eso convierte la matriz en un \texttt{Map} que se
pasa como acumulador de un \texttt{reduce}, sin perder nada de eficiencia y
sin trucos.

\paragraph{El reloj del sistema no alcanza para medir.}
El enunciado pide reportar el tiempo de validación. La primera versión usaba
\texttt{:timer.tc} y reportaba \texttt{0.000 ms}. Investigando resultó que el
reloj monotónico del sistema operativo tiene una resolución de unos 100
microsegundos en Windows, y una validación tarda unos 14. Es decir, el
instrumento era más grueso que lo que se quería medir. La solución fue repetir
la validación hasta superar un umbral de tiempo y dividir el total entre el
número de corridas.

\paragraph{El calentamiento de la máquina virtual.}
Relacionado con lo anterior: la primera ejecución medía 1.3 ms y las
siguientes 0.013 ms, cien veces menos. La causa es el compilador JIT de la
BEAM, que optimiza el código recién después de ejecutarlo varias veces.
Reportar la primera corrida habría dado un número engañoso. El esquema de
promediar varias corridas resuelve también este problema.

\paragraph{Variables duplicadas en la binarización.}
La primera versión de la binarización creaba una variable nueva por cada
cuerpo largo que encontraba. Como \gram{( E )} aparece en tres reglas
distintas, generaba \gram{X1}, \gram{X3} y \gram{X5} con el mismo contenido
\gram{E T2}, dando 22 producciones donde el CNF de referencia tiene 20. Se
agregó un registro de pares ya encadenados para reutilizarlos. Con eso el
resultado coincide exactamente en tamaño y estructura con el del curso.

\paragraph{La gramática del enunciado ya estaba en CNF.}
Al correr los cuatro pasos sobre la gramática del proyecto, ninguno cambió
nada. No era un error: revisando regla por regla, todas son ya de la forma
$A \to BC$ o $A \to a$, no hay épsilon, no hay unitarias y no hay símbolos
inútiles. \gram{VP -> cooks} parece unitaria pero no lo es, porque
\gram{cooks} es terminal, no variable. Por eso agregamos la gramática
aritmética del curso, que sí ejercita los cuatro pasos de principio a fin.

\paragraph{Una ventana que colgaba el programa.}
El árbol se abre en una ventana aparte. La primera implementación usaba
\texttt{cmd /c start}, y el programa se quedaba congelado hasta que uno
cerrara el navegador. La razón es que el navegador heredaba la salida estándar
del proceso y Elixir esperaba a que se cerrara. Se cambió a
\texttt{explorer.exe}, que lanza el visor y devuelve de inmediato.

\paragraph{Una prueba mal escrita.}
Habíamos puesto \textit{the cake eats} en la lista de frases que deberían
rechazarse, y la prueba falló. Revisando, el programa tenía razón y nosotros
no: la frase es sintácticamente correcta. Pasó a la lista de aceptadas, y
terminó siendo el mejor ejemplo de la advertencia del enunciado.

\subsection{Recomendaciones}

\begin{itemize}[leftmargin=*]
  \item \textbf{Separar la lógica de la interfaz desde el principio.} Permitió
  escribir pruebas sin simular una terminal, y cambiar la presentación sin
  riesgo de romper los algoritmos.
  \item \textbf{Usar el ejemplo resuelto del curso como prueba automática.}
  Tener \texttt{1.txt} y \texttt{1-cnf.txt} y poder verificar la equivalencia
  entre ambos dio mucha más confianza que revisar la salida a ojo.
  \item \textbf{Desconfiar de las mediciones de tiempo muy chicas.} Si el
  número sale 0 o varía cien veces entre corridas, el problema casi siempre
  está en el instrumento, no en el algoritmo.
  \item \textbf{Cuando una prueba falla, revisar primero la prueba.} En nuestro
  caso el código estaba bien y el error era nuestro entendimiento de la
  gramática.
  \item \textbf{Para una próxima versión}, se podría mostrar todos los árboles
  cuando la frase es ambigua, en vez de uno solo más la cuenta. La información
  ya está en los backpointers; solo falta recorrerlos exhaustivamente.
\end{itemize}

% ====================================================================
\section{Conclusiones}

El programa cumple lo que pide el enunciado: implementa el algoritmo de
simplificación de gramáticas, implementa CYK, responde sí o no, reporta el
tiempo de validación y construye el árbol de derivación.

Más allá de eso, el ejercicio dejó claro por qué la Forma Normal de Chomsky no
es un requisito arbitrario: es la condición que hace posible atacar el problema
con programación dinámica. Sin la restricción de dos símbolos por cuerpo, no
hay dónde partir la subcadena y el algoritmo no existe.

También quedó claro que la elección de un lenguaje funcional, que al principio
parecía un obstáculo para implementar programación dinámica, terminó siendo
una ventaja. Los tres pasos de simplificación son cálculos de punto fijo y se
escriben casi igual que en la pizarra, y la tabla de CYK resultó no necesitar
mutación en ningún momento.

% ====================================================================
\section{Referencias}

\begin{itemize}[leftmargin=*]
  \item Hopcroft, J., Motwani, R. y Ullman, J. \textit{Introduction to Automata
  Theory, Languages, and Computation}. 3ª edición.
  \item Aho, A., Lam, M., Sethi, R. y Ullman, J. \textit{Compiladores:
  principios, técnicas y herramientas}. 2ª edición.
  \item Material del curso: presentación de introducción a gramáticas y laboratorios.
\end{itemize}

\end{document}
